\section{Introduction}
\label{sec:introduction}

\subsection{Context}

The Trusted Agents Protocol (TAP) is an open-source, local-first protocol stack that lets
personal AI agents---an OpenClaw or Hermes assistant run by one person---find, verify and
message the agents run by people they already know~\cite{tap-repo}. Each agent holds an
on-chain identity as an ERC-8004 token~\cite{erc8004}, publishes a registration file
describing its messaging endpoint, connects to peers through signed invitation links, and
exchanges JSON-RPC~2.0 payloads over the XMTP messaging network~\cite{xmtp}. Owners
control what a peer may ask for through \emph{directional permission grants}, for example
``may request up to 10~USDC per week''. TAP is explicitly built for trust between known
humans mediated through their agents, not for an open marketplace.

Taiko, the industry partner for this capstone, operates an Ethereum-based rollup
(Taiko Alethia) and had already integrated with TAP at the identity layer: agents can
register on Taiko (\code{eip155:167000}) and can upload their registration file through
\emph{Tack} (\code{tack.taiko.xyz}), Taiko's x402-metered IPFS upload
service~\cite{x402}. The project brief extended this from paid storage to a broader notion
of paid interaction between agents.

\subsection{Problem}

Once two agents are connected, TAP as of April~2026 delivered every peer message into the
receiving agent's context without any notion of cost. This creates three concrete
problems, all observed in the codebase before this work:

\begin{enumerate}
  \item \textbf{Unmetered attention.} A peer message costs the \emph{receiver} LLM tokens on
        every prompt into which it is rendered, and in streaming hosts it can wake an idle
        agent. The sender pays nothing. This is structurally the same externality that
        made e-mail spam economically rational~\cite{dwork-naor,hashcash}.
  \item \textbf{Buried escalations.} The notification block injected into the agent's prompt
        was rendered first-in-first-out with a hard cap of twenty lines. A burst of chatter
        arriving before a transfer-approval request pushed the request into an opaque
        ``$N$ more omitted'' footer. The one line the operator most needed to see was the
        one most likely to be hidden.
  \item \textbf{No vocabulary for pricing.} Even if a receiver wanted to charge for attention,
        the protocol had no way to advertise a price, no error that a sender could interpret
        as ``payment required'', and no unit in which attention could be measured or billed.
\end{enumerate}

\subsection{Objectives}

The capstone set out to make LLM attention a first-class, measurable and optionally
prepaid resource on top of TAP's existing invite-and-grant trust baseline. Concretely:

\begin{enumerate}[label=O\arabic*.]
  \item Reduce the context a message flood injects into the receiver's prompt, and
        guarantee that escalations are always rendered.
  \item Measure attention spend per peer so that pricing decisions rest on data.
  \item Let an agent advertise attention prices and reject un-granted, unpaid messages
        with a machine-readable quote, without changing the behaviour of any agent that
        does not opt in.
  \item Provide prepaid postage so that a sender pays once and spends many times, with
        replay and double-spend protection.
  \item Let a sender pay for an escalation (a ``priority'' message) and let a receiver
        meter the free tier of a grant holder through a weekly notification quota.
  \item Validate the mechanisms with a reproducible evaluation.
\end{enumerate}

\subsection{A change of approach, stated up front}

The first progress journal proposed to meet these objectives with a standalone
TypeScript service modelled on Tack: senders would pay per message through x402 on Taiko
Alethia, and the service would gate delivery. The implemented system does not do this.
Section~\ref{sec:design-pivot} sets out why: the scarce resource is the receiver's
attention, which only the receiver's own daemon can observe and account for; TAP already
supplies authenticated peer identity and grants that a standalone service would have to
duplicate; and the upstream TAP maintainers' research note for the messaging layer argued for keeping
chain RPC out of the hot path in favour of prepaid credits~\cite{tap-research-note}. The
change was a design decision taken with the full time budget available, not a fallback.
The x402 idea survives in the implemented system as a ``402 over XMTP'': a JSON-RPC
\code{-32050} error whose \code{data} carries the same price list the agent publishes on
chain.

\subsection{Contributions and division of work}

This is an individual report on a two-person project. The author (Yiqun Xu) designed and
implemented the entire paid-attention layer described in
Sections~\ref{sec:design}--\ref{sec:implementation}: 23~implementation commits on the
branch \path{feat/phase1-notification-coalescing}, roughly 8{,}200 added lines across
91~files, including a new workspace package \code{@trustedagents/app-postage}, mock
end-to-end test phases~6--8, and the accompanying documentation in the TAP skill file. The
teammate (Zhang Hanyu) led the earlier design exploration and owned the evaluation:
the benchmark scripts and \path{evaluation/results.md} merged as pull request~\#2. All
figures in Section~\ref{sec:evaluation} are his and are reproduced exactly.
Section~\ref{sec:project-management} gives the full account, including the fact that the
industry mentor has been unreachable since late August~2026 and that the final design
therefore carries no mentor sign-off.

\subsection{Report structure}

Section~\ref{sec:background} introduces TAP, the surrounding standards and the related
work on attention pricing. Section~\ref{sec:requirements} states the requirements and the
threat model. Section~\ref{sec:design} presents the design and its rationale, including
the pivot away from a standalone service. Section~\ref{sec:implementation} describes the
implementation. Section~\ref{sec:evaluation} reports the evaluation.
Section~\ref{sec:discussion} discusses limitations and security considerations.
Section~\ref{sec:project-management} covers project management and the division of
labour. Section~\ref{sec:conclusion} concludes with future work. Appendices reproduce
the operator-facing configuration, the wire formats and the commit log.
